\section{Сборка и проверка работы}
\subsection{Среда сборки}
Сборка выполнена в Windows из PowerShell. В журнале CMake определены генератор Visual Studio 18 2026, компилятор MSVC 19.51.36260.0 и Windows SDK 10.0.26100.0. Использована конфигурация Debug.

Из корня внешнего репозитория \texttt{OS\_Labs} выполнены команды:
\begin{verbatim}
cmake -S .\Operation-Systems -B .\Operation-Systems\build
cmake --build .\Operation-Systems\build --config Debug
\end{verbatim}

Сборка завершилась без ошибок компиляции и компоновки. Созданы библиотеки \texttt{os.lib}, \texttt{lab2\_logic.lib} и исполняемые файлы \texttt{lab2\_parent.exe}, \texttt{lab2\_child.exe}.

\subsection{Запуск программы}
Родительская программа запускается командой:
\begin{verbatim}
.\Operation-Systems\build\bin\lab2\lab2_parent.exe
\end{verbatim}

Обе программы должны находиться в одной папке. Имя дочернего исполняемого файла в функции определения пути должно совпадать с именем результата сборки. На этапе проверки обнаружено несовпадение имён \texttt{child.exe} и \texttt{lab2\_child.exe}: оно приводило к ошибке \texttt{CreateProcessW}. После согласования имени запуск был подтверждён.

\subsection{Контрольный пример}
Для проверки основного сценария используется следующая последовательность ввода:
\begin{verbatim}
result.txt
4
6
9
17
\end{verbatim}

Ожидаемое содержимое файла:
\begin{verbatim}
4
6
9
\end{verbatim}

Значения 4, 6 и 9 составные и подлежат записи. Значение 17 простое: оно не записывается и завершает обработку. При запуске из корня \texttt{OS\_Labs} относительный путь \texttt{result.txt} относится к этой рабочей папке, а не к папке исполняемого файла.

\subsection{Дополнительные контрольные случаи}
В таблице приведены ожидаемые результаты для повторной проверки. Она описывает набор проверок, а не журнал их выполнения на пользовательском компьютере.

\begin{center}
\begin{tabular}{p{0.25\textwidth}p{0.66\textwidth}}
Входные данные & Ожидаемый результат \\
\hline
\texttt{8, -4} & Запись 8, завершение на отрицательном значении \\
\texttt{2} & Завершение без записи числа \\
\texttt{0, 1, 4, -1} & Пропуск 0 и 1, запись 4, завершение на -1 \\
\texttt{abc, 4, -1} & Сообщение об ошибке ввода, запись 4, завершение \\
Два числа в строке & Отклонение команды без завершения обмена \\
Выход за диапазон \texttt{int} & Отклонение команды \\
Конец ввода & Закрытие входного канала и ожидание ребёнка \\
Недоступный выходной путь & Диагностика ошибки и ненулевой код завершения \\
\end{tabular}
\end{center}

В данной работе исследуется корректность межпроцессного обмена. Измерения ускорения и эффективности многопоточного алгоритма не проводились: они не относятся к заданию лабораторной работы №2. Подтверждение сборки Windows не рассматривается как подтверждение запуска Linux-реализации на том же компьютере.
